% Part: lambda-calculus
% Chapter: introduction
% Section: reduction

\documentclass[../../../include/open-logic-section]{subfiles}

\begin{document}

\olfileid[hi]{lam}{int}{red}
\olsection{लैम्ब्डा पदों का अपचयन}

लैम्ब्डा पदों के साथ क्या किया जा सकता है? उन्हें सरल किया जा सकता है। यदि
$M$ और $N$ कोई भी लैम्ब्डा पद और $x$ कोई भी चर हो, तो~$M$ में~$x$ के
स्थान पर $N$ प्रतिस्थापित करने से मिलने वाले परिणाम को हम
$\Subst{M}{N}{x}$ से दर्शा सकते हैं; इससे पहले~$M$ के उन आबद्ध चरों का
पुनर्नामकरण कर दिया जाता है जो प्रतिस्थापन के बाद~$N$ के मुक्त चरों से
टकराते। उदाहरणार्थ,
\[
\Subst{(\lambd[w][xxw])}{yyz}{x} = \lambd[w][(yyz)(yyz)w].
\]

\begin{digress}
प्रतिस्थापन की वैकल्पिक संकेत-लिपियाँ $[N/x]M$, $[x/N]M$ और $M[x/N]$ भी
हैं। सावधान रहें!
\end{digress}

सहज रूप से $(\lambd[x][M])N$ और $\Subst{M}{N}{x}$ का अर्थ एक ही है; पहले
पद को दूसरे से बदलने की क्रिया को \emph{$\beta$-संकुचन} कहते हैं।
$(\lambd[x][M])N$ को \emph{रिडेक्स} और $\Subst{M}{N}{x}$ को उसका
\emph{संकुचित रूप} कहते हैं। सामान्यतः, यदि किसी उपपद के $\beta$-संकुचन
द्वारा पद $P$ को~$P'$ में बदला जा सकता है, तो हम कहते हैं कि $P$
\emph{एक चरण में $\beta$-अपचयित होकर $P'$ बनता है}, और $P \redone P'$
लिखते हैं। यदि $P$ से एक-चरणीय अपचयनों की किसी संख्या (शून्य भी संभव है)
द्वारा~$P'$ प्राप्त किया जा सकता है, तो $P$ \emph{$\beta$-अपचयित होकर}
~$P'$ बनता है; संकेतों में, $P \red P'$। जिस पद का आगे $\beta$-अपचयन न
किया जा सके वह \emph{$\beta$-अनपचयनीय}, अथवा \emph{$\beta$-प्रसामान्य}
कहलाता है। प्रसंग स्पष्ट होने पर हम ``$\beta$-अपचयित होता है'' आदि के बदले
केवल ``अपचयित होता है'' कहेंगे।

कुछ उदाहरणों पर विचार करें।
\begin{enumerate}
\item हमारे पास है
\begin{align*}
(\lambd[x][xxy]) \lambd[z][z] & \redone (\lambd[z][z])(\lambd[z][z]) y \\
& \redone (\lambd[z][z]) y \\
& \redone y.
\end{align*}
\item किसी पद को ``सरल करना'' उसे अधिक जटिल बना सकता है:
\begin{align*}
(\lambd[x][xxy])(\lambd[x][xxy]) & \redone (\lambd[x][xxy])(\lambd[x][xxy])y \\
& \redone (\lambd[x][xxy])(\lambd[x][xxy])yy \\
& \redone \dots
\end{align*}
\item यह किसी पद को अपरिवर्तित भी छोड़ सकता है:
\[
(\lambd[x][xx])(\lambd[x][xx]) \redone (\lambd[x][xx])(\lambd[x][xx]).
\]
\item कुछ पदों को एक से अधिक प्रकार से अपचयित भी किया जा सकता है;
  उदाहरणार्थ,
\[
(\lambd[x][(\lambd[y][yx]) z)] v \redone (\lambd[y][yv]) z
\]
सबसे बाहरी अनुप्रयोग को संकुचित करने पर; और
\[
(\lambd[x][(\lambd[y][yx]) z)] v \redone (\lambd[x][zx]) v
\]
सबसे भीतरी अनुप्रयोग को संकुचित करने पर। फिर भी ध्यान दें कि इस स्थिति में
दोनों पद आगे अपचयित होकर एक ही पद~$zv$ बनते हैं।
\end{enumerate}

अंतिम उदाहरण का अंतिम परिणाम संयोग नहीं है; वह लैम्ब्डा कलन के एक गहरे और
महत्त्वपूर्ण गुण को दर्शाता है, जिसे ``चर्च--रॉसर गुण'' कहते हैं।

\end{document}
